
\SGAInlineRefTarget{sga3:VIIA:nde:2}\section*{2. Invariant Differential Operators on Group Schemes}
\phantomsection
\label{sga3:VIIA:2}
\addcontentsline{toc}{section}{2. Invariant Differential Operators on Group Schemes}

\medskip
\noindent\textbf{2.1.}
\phantomsection
\label{sga3:VIIA:2.1}
Let \(G\) be an \(S\)-group scheme; we denote its unit section by
\(\varepsilon\), or by \(\varepsilon_G:S\to G\).

\medskip
\noindent\textbf{Definition.}
\textit{Let \(U(G)\) be the \(\Gamma(\mathcal O_S)\)-module of
\(S\)-deviations of \(\varepsilon_G\) (or \(S\)-deviations of the
origin), as in \SGARefLink{sga3:VIIA:1.1}{1.1}.}

\textit{If \(d,e\in U(G)\), then \(d\times e\) is an \(S\)-deviation of
\[
  \varepsilon\times\varepsilon:
  S\simeq S\times S\longrightarrow G\times G.
\]
The image of \(d\times e\) under the multiplication morphism
\(m:G\times G\to G\), as in \SGARefLink{sga3:VIIA:1.2}{1.2}, is called the \emph{product} of
\(d\) and \(e\) and is denoted \(d\cdot e\).}

\textit{The \(\Gamma(\mathcal O_S)\)-module \(U(G)\) is thereby made
into an associative \(\Gamma(\mathcal O_S)\)-algebra with unit
\(\varepsilon_G\). We call \(U(G)\) the \emph{infinitesimal algebra}
of \(G\).}\footnote{This is now usually called the algebra of
distributions (at the origin) of \(G\); see Demazure--Gabriel,
II.4, \SGARefLink{sga3:VIIA:6.1}{6.1}, and Jantzen, I.7.7.}

As \(T\) ranges over schemes above \(S\), the infinitesimal algebra
\(U(G_T)\) of the \(T\)-group \(G\times T\) plainly varies
contravariantly with \(T\); thus one may speak of the
\emph{infinitesimal-algebra functor}.

Restricting \(T\) to the open subsets of \(S\) gives a presheaf
\(T\mapsto U(G_T)\) of \(\mathcal O_S\)-algebras. By \SGARefLink{sga3:VIIA:1.1.3}{1.1.3} it is a
sheaf. We denote it by \(\mathcal U(G)\) and call it the
\emph{sheaf of infinitesimal algebras of \(G\)}.

The algebra \(U(G)\) is also covariantly functorial in \(G\). Indeed,
if \(u:G\to H\) is a homomorphism of \(S\)-groups and \(d\) is an
\(S\)-deviation of \(\varepsilon_G\), its image under \(u\) is the
element
\[
  U(u)(d)=u d
\]
of \(U(H)\). The resulting map \(U(u):U(G)\to U(H)\) is a
homomorphism of \(\Gamma(\mathcal O_S)\)-algebras. Likewise one obtains
a homomorphism
\(\mathcal U(u):\mathcal U(G)\to\mathcal U(H)\).

\medskip
\noindent\textbf{2.2.}
\phantomsection
\label{sga3:VIIA:2.2}
Let \(d\in U(G)\), that is, let \(d\) be an \(S\)-deviation of the
origin of \(G\). The base change of \(d\) is the \(S\)-deviation
\[
  d\times G
\]
of
\(\varepsilon\times G:G\simeq S\times G\to G\times G\), as in
\SGARefLink{sga3:VIIA:1.2.2}{1.2.2}. Its image under multiplication \(m:G\times G\to G\) is an
\(S\)-deviation of
\[
  m\circ(\varepsilon\times\operatorname{id}_G)
  =\operatorname{id}_G,
\]
hence an element of \(\operatorname{Dif}_{G/S}\), which we denote
by \(d^G\).

The map \(d\mapsto d^G\) is plainly
\(\Gamma(\mathcal O_S)\)-linear. The following commutative diagram
shows that
\[
  (e\cdot d)^G=d^G\cdot e^G.
\]
\begin{center}
  \phantomsection
  \label{sga3:VIIA:2.2-right-translation-associativity}
  \begin{tikzcd}[column sep=small,row sep=large]
&&
G\times G\times G
  \arrow[rr,"m\times G"]
  \arrow[dr,"G\times m"]
&&
G\times G
  \arrow[dd,"m"]
\\
&
G\times G
  \arrow[ur,"G\times d\times G"]
  \arrow[dr,"m"]
&&
G\times G
  \arrow[dr,"m"]
\\
G
  \arrow[ur,shift left=.7ex,"e\times G"]
  \arrow[ur,shift right=.7ex,"\varepsilon\times G"']
  \arrow[rr,shift left=.55ex,"e^G"]
  \arrow[rr,shift right=.55ex,"\operatorname{id}_G"']
&&
G
  \arrow[ur,shift left=.7ex,"d\times G"]
  \arrow[ur,shift right=.7ex,"\varepsilon\times G"']
  \arrow[rr,shift left=.55ex,"d^G"]
  \arrow[rr,shift right=.55ex,"\operatorname{id}_G"']
&&
G
\end{tikzcd}

\end{center}
\footnote{The editors corrected the original diagram by replacing
\(d\times G\times G\) with \(G\times d\times G\). The composite on
the left side of the upper triangle is then
\((e\times d)\times G\), and \(d\mapsto d^G\) becomes an
anti-isomorphism from \(U(G)\) to the right-invariant differential
operators (\SGARefLink{sga3:VIIA:2.3}{2.3}--\SGARefLink{sga3:VIIA:2.4}{2.4}). If instead \({}^Gd\) is defined as the image of
\(G\times d\) under \(m\), one similarly obtains an isomorphism onto
the left-invariant differential operators. \SGARefLink{sga3:VIIA:2.4}{Sections 2.4} and \SGARefLink{sga3:VIIA:2.5}{2.5} have
been corrected accordingly.}

The two lower triangles commute by the definitions of \(d^G\) and
\(e^G\). The composite \(S\)-deviation of \(e\times G\) and
\(G\times d\times G\) is \((e\times d)\times G\), by \SGARefLink{sga3:VIIA:1.2.2}{1.2.2}. Its image
under \(m\times G\) is \((e\cdot d)\times G\), and the image of the
latter under \(m\) is therefore \((e\cdot d)^G\).

Thus one obtains an antihomomorphism
\[
  U(G)\longrightarrow\operatorname{Dif}_{G/S}
\]
of \(\Gamma(\mathcal O_S)\)-algebras, called \emph{right
translation}.\footnote{It would be preferable to call this a
\emph{left action}. For example, if \(d\) is an \(S\)-derivation of
the origin, then by \SGARefLink{sga3:VIIA:1.2.1}{1.2.1} it is the composite of
\((\tau,\partial_t):S\to I_S\) with a morphism \(x:I_S\to G\)
satisfying \(x\circ\tau=\varepsilon\). The derivation \(d^G\) of
\(\mathcal O_G\) sends a local section \(\varphi\) to the section
\(g\mapsto\partial_t\varphi(xg)\). With this terminology one could say
that the left action commutes with right translations.}

If \(\mathcal{Dif}_{G/S}\) denotes the sheaf of \(S\)-differential
operators on \(G\), as in \SGARefLink{sga3:VIIA:1.4}{1.4}, and \(p:G\to S\) is the structural
morphism, one likewise defines a right translation
\[
  \mathcal U(G)\longrightarrow p_*(\mathcal{Dif}_{G/S}).
\]

\medskip
\noindent\textbf{2.3.}
\phantomsection
\label{sga3:VIIA:2.3}
We now characterize the differential operators on \(G/S\) of the
form \(d^G\). Let \(g:S\to G\) be a section of the structural
morphism and let \(g_G\) denote right translation by \(g\), namely
\[
  g_G:
  G\simeq G\times S
  \xrightarrow{G\times g}G\times G
  \xrightarrow{m}G.
\]
For every differential operator \(D\) on \(G/S\), the composite
\[
  g_G^{-1}Dg_G
\]
is again an \(S\)-deviation of \(\operatorname{id}_G\), hence an
element of \(\operatorname{Dif}_{G/S}\). We write
\[
  D^g=g_G^{-1}Dg_G.
\]
We call \(D\) \emph{right-invariant} if for every base change
\(t:T\to S\) and every section \(g:T\to G\times T\), one has
\[
  (D_T)^g=D_T.
\]

\medskip
\noindent\textbf{Lemma.}
\phantomsection
\label{sga3:VIIA:2.3-invariance-lemma}
\textit{For every differential operator \(D\) on \(G/S\), the
following assertions are equivalent, where \(m\) is the
multiplication morphism of \(G\):}
\begin{enumerate}
\item[\textit{(i)}] \textit{\(D\) is right-invariant.}
\item[\textit{(ii)}] \textit{The following two deviations of \(m\)
are equal:}
\[
  Dm=m(D\times G).
\]
\end{enumerate}

\noindent\((ii)\Rightarrow(i)\).
Since condition \SGARefLink{sga3:vii:paragraph:5-1:equation:ii}{(ii)} is stable under base change, it suffices to show
that it implies \(D^g=D\) for every section \(g:S\to G\). Let
\[
  h=G\times g:G\simeq G\times S\longrightarrow G\times G,
\]
so that \(m\circ h=g_G\). The equality \(D^g=D\) is equivalent to
\(g_G\circ D=D\circ g_G\), and this follows from the commutative
diagram
\begin{center}
  \phantomsection
  \label{sga3:VIIA:2.3-ii-implies-i}
  {\small
\begin{tikzcd}[column sep=large,row sep=large]
G
  \arrow[d,shift right=.55ex,"D"']
  \arrow[d,shift left=.55ex,"\operatorname{id}_G"]
&
G\times G
  \arrow[l,"m"']
  \arrow[d,shift right=.55ex,"D\times G"']
  \arrow[d,shift left=.55ex,"\operatorname{id}_{(G\times G)}"]
&
G\,.
  \arrow[l,"h"']
  \arrow[d,shift right=.55ex,"D"']
  \arrow[d,shift left=.55ex,"\operatorname{id}_G"]
\\
G
&
G\times G
  \arrow[l,"m"']
&
G
  \arrow[l,"h"']
\end{tikzcd}}

\end{center}

\noindent\((i)\Rightarrow(ii)\).
Take \(t:T\to S\) to be the structural morphism \(p:G\to S\), and
take the section \(g:T\to G\times T\) to be the diagonal
\(\Delta:G\to G\times G\). Right translation
\[
  \Delta_{G\times G}:G\times G\longrightarrow G\times G
\]
then has components \(m\) and \(\operatorname{pr}_2\). The equality
\((D_G)^\Delta=D_G\) is equivalent to commutativity of the first
square in
\begin{center}
  \phantomsection
  \label{sga3:VIIA:2.3-i-implies-ii}
  {\small
\begin{tikzcd}[column sep=large,row sep=large]
G\times G
  \arrow[r,"\Delta_{G\times G}"]
  \arrow[d,shift right=.55ex,"D_G"']
  \arrow[d,shift left=.55ex,"\operatorname{id}_{G\times G}"]
&
G\times G
  \arrow[r,"\operatorname{pr}_1"]
  \arrow[d,shift right=.55ex,"D_G"']
  \arrow[d,shift left=.55ex,"\operatorname{id}_{G\times G}"]
&
G\,.
  \arrow[d,shift right=.55ex,"D"']
  \arrow[d,shift left=.55ex,"\operatorname{id}_G"]
\\
G\times G
  \arrow[r,"\Delta_{G\times G}"]
&
G\times G
  \arrow[r,"\operatorname{pr}_1"]
&
G
\end{tikzcd}}

\end{center}
because \(m=\operatorname{pr}_1\circ\Delta_{G\times G}\). This proves
\SGARefLink{sga3:vii:paragraph:5-1:equation:ii}{(ii)}.

Now consider \(d\in U(G)\). Every square in the following diagram is
commutative:
\begin{center}
  \phantomsection
  \label{sga3:VIIA:2.3-origin-deviation-invariant}
  {\scriptsize
\begin{tikzcd}[column sep=small,row sep=large]
G\times G
  \arrow[r,equal]
  \arrow[d,"m"']
&
S\times G\times G
  \arrow[r,shift left=.55ex,"d\times G\times G"]
  \arrow[r,shift right=.55ex,"\varepsilon\times G\times G"']
  \arrow[d,"S\times m"']
&
G\times G\times G
  \arrow[r,"m\times G"]
  \arrow[d,"G\times m"]
&
G\times G
  \arrow[d,"m"]
\\
G
  \arrow[r,equal]
&
S\times G
  \arrow[r,shift left=.55ex,"d\times G"]
  \arrow[r,shift right=.55ex,"\varepsilon\times G"']
&
G\times G
  \arrow[r,"m"]
&
G
\end{tikzcd}}

\end{center}
Since
\[
  m\circ(d\times G)=d^G
  \qquad\text{and}\qquad
  (m\times G)\circ(d\times G\times G)=d^G\times G,
\]
one also has
\[
  d^G\circ m=m\circ(d^G\times G).
\]
Thus, for every \(S\)-deviation \(d\) of the origin, \(d^G\) is a
right-invariant differential operator.

\medskip
\noindent\textbf{Theorem 2.4.}
\phantomsection
\label{sga3:VIIA:2.4}
\begin{enumerate}
\item[\textit{\SGAInlineRefTarget{sga3:VIIA:2.4:item:i}(i)}]
\textit{The map \(d\mapsto d^G\) is an anti-isomorphism of the
infinitesimal algebra \(U(G)\) onto the subalgebra
\(\operatorname{Dif}^{G}_{G/S}\) of \(\operatorname{Dif}_{G/S}\)
formed by the right-invariant differential operators.}
\item[\textit{\SGAInlineRefTarget{sga3:VIIA:2.4:item:ii}(ii)}]
\textit{Likewise, the map \(d\mapsto{}^Gd\) is an isomorphism from
\(U(G)\) onto the subalgebra of \(\operatorname{Dif}_{G/S}\) formed
by the left-invariant differential operators.}
\end{enumerate}
\footnote{The editors corrected ``isomorphism'' to
``anti-isomorphism'' and added assertion \SGARefLink{sga3:vii:paragraph:5-1:equation:ii}{(ii)}; see the preceding
editorial note.}

Let \(D\) be an arbitrary differential operator on \(G/S\), and let
\(D_0\) denote its \emph{value at the origin}, that is, the composite
deviation
\[
  S\xrightarrow{\varepsilon}G
   \mathop{\longrightarrow}^{D}_{\operatorname{id}_G}G.
\]
The right-invariant differential operator \((D_0)^G\) is obtained by
the composite
\begin{center}
  \phantomsection
  \label{sga3:VIIA:2.4-D0-composite}
  {\small
\begin{tikzcd}[column sep=normal]
G\simeq S\times G
  \arrow[r,"\varepsilon\times G"]
&
G\times G
  \arrow[r,shift left=.55ex,"D\times G"]
  \arrow[r,shift right=.55ex,"\operatorname{id}_{G\times G}"']
&
G\times G
  \arrow[r,"m"]
&
G\,.
\end{tikzcd}}

\end{center}
If \(D\) is right-invariant, then \(Dm=m(D\times G)\), and hence
\[
  D=Dm(\varepsilon\times G)
   =m(D\times G)(\varepsilon\times G)
   =(D_0)^G.
\]
Thus \(d\mapsto d^G\) is surjective.

Conversely, for every \(S\)-deviation \(d\) of the origin one has the
commutative square
\begin{center}
  \phantomsection
  \label{sga3:VIIA:2.4-origin-value-square}
  \begin{tikzcd}[column sep=huge,row sep=large]
G\times G
&
G
  \arrow[l,"d\times G"']
\\
G\times S\simeq G
  \arrow[u,"G\times\varepsilon"]
&
S
  \arrow[l,"d"']
  \arrow[u,"\varepsilon"']
\end{tikzcd}

\end{center}
from which
\[
  d=m(G\times\varepsilon)d
   =m(d\times G)\varepsilon
   =(d^G)_0.
\]
\textit{A fortiori}, \(d\mapsto d^G\) is injective. This proves the
theorem.

As \(S\) varies, \SGARefLink{sga3:VIIA:2.4}{Theorem~2.4} also shows that right translation
\[
  \mathcal U(G)\longrightarrow p_*(\mathcal{Dif}_{G/S})
\]
is an anti-isomorphism from \(\mathcal U(G)\) onto the sheaf of
\(\mathcal O_S\)-algebras
\[
  p_*(\mathcal{Dif}_{G/S})^G,
\]
which assigns to every open \(V\subset S\) the algebra
\(\operatorname{Dif}^{G_V}_{G_V/V}\).

\medskip
\noindent\textbf{Remark 2.4.1.}
\phantomsection
\label{sga3:VIIA:2.4.1}
Consider the commutative diagram
\begin{center}
  \phantomsection
  \label{sga3:VIIA:2.4.1-eta-cartesian-square}
  \begin{tikzcd}[column sep=huge,row sep=large]
G
  \arrow[d,shift right=.55ex,"p"']
&
G\times G
  \arrow[l,"\eta"']
  \arrow[d,shift right=.55ex,"\operatorname{pr}_1"']
\\
S
  \arrow[u,shift right=.55ex,"\varepsilon"']
&
G,
  \arrow[l,"p"']
  \arrow[u,shift right=.55ex,"\Delta"']
\end{tikzcd}

\end{center}
where \(\eta\) denotes the morphism
\((x,y)\mapsto yx^{-1}\).\footnote{Thus \(G\) acts on itself on the
left by right translations.}
It induces morphisms
\[
  \eta':\eta^{-1}(\mathcal O_G)\longrightarrow\mathcal O_{G\times G}
\]
and
\[
  \Delta^{-1}(\eta'):
  p^{-1}\varepsilon^{-1}(\mathcal O_G)
  \longrightarrow\Delta^{-1}(\mathcal O_{G\times G}).
\]

For every integer \(n\geq1\), set
\[
  \mathfrak p^n_{G/S}
  =
  \varepsilon^{-1}(\mathcal O_G)/
  \mathcal I_\varepsilon^{\,n+1}.
\]
\footnote{The editors corrected the original here: the relevant
cartesian square is formed by \(\varepsilon,\eta,\Delta,p\), not by
\(p,p,\eta,\operatorname{pr}_1\).}
Since the square formed by \(\varepsilon,\eta,\Delta,p\) is cartesian,
\(\Delta^{-1}(\eta')\) induces an isomorphism of
\(\mathcal O_G\)-modules
\[
  p^*(\mathfrak p^n_{G/S})
  \xrightarrow{\sim}
  \mathcal P^n_{G/S}.
\]
The differential operators on \(G/S\) of order \(\leq n\) therefore
correspond bijectively to the homomorphisms of
\(\mathcal O_G\)-modules
\[
  p^*(\mathfrak p^n_{G/S})\longrightarrow\mathcal O_G,
\]
or equivalently to the homomorphisms of \(\mathcal O_S\)-modules
\[
  \mathfrak p^n_{G/S}\longrightarrow p_*(\mathcal O_G).
\]
Under this bijection the right-invariant differential operators
correspond to the composites
\begin{center}
  \phantomsection
  \label{sga3:VIIA:2.4.1-invariant-factorization}
  \(
\mathfrak p^n_{G/S}
  \longrightarrow
\mathcal O_S
  \xrightarrow{\mathrm{can.}}
p_*(\mathcal O_G)\,.
\)

\end{center}
This recovers the isomorphism of \SGARefLink{sga3:VIIA:2.4}{Theorem~2.4}.

\medskip
\noindent\textbf{2.5.}
\phantomsection
\label{sga3:VIIA:2.5}
\footnote{The editors reordered this paragraph, beginning with the
definition of \(\alpha:\operatorname{Lie}(G)\to U(G)\), and corrected
the original as indicated in the note to \SGARefLink{sga3:VIIA:2.2}{2.2}.}
Let \(\operatorname{Lie}(G)\) be the Lie algebra of \(G\).\footnote{In
this expos\'e, for an \(S\)-group scheme \(G\) (respectively, an
\(S\)-scheme \(X\)), \(\operatorname{Lie}(G)\) (respectively,
\(\operatorname{Lie}(\operatorname{Aut}X)\)) denotes, in the notation
of Expos\'e II,
\(\operatorname{Lie}(G/S)(S)\) (respectively,
\(\operatorname{Lie}(\operatorname{Aut}_S(X)/S)(S)\)). It is a
\(\Gamma(\mathcal O_S)\)-Lie algebra by II, \SGARefLink{sga3:ii:subsection:4-11}{4.11} and \SGARefLink{sga3:ii:subsection:3-14}{3.14}.}
We shall define a homomorphism of
\(\Gamma(\mathcal O_S)\)-Lie algebras
\[
  \alpha:\operatorname{Lie}(G)\longrightarrow U(G).
\]

Let \(s:S\to I_S\) be the zero section and let \(\sigma\) be the
deviation of \(s\) defined in \SGARefLink{sga3:VIIA:1.2.0}{1.2.0}. Recall from II, \SGARefLink{sga3:ii:subsection:4-1}{4.1} that
\(\operatorname{Lie}(G)\) is the set of morphisms \(x:I_S\to G\)
such that \(x\circ s=\varepsilon_G\). The composite
\begin{center}
  \phantomsection
  \label{sga3:VIIA:2.5-sigma-x}
  \begin{tikzcd}[column sep=large]
S
  \arrow[r,shift left=.55ex,"\sigma"]
  \arrow[r,shift right=.55ex,"s"']
&
I_S
  \arrow[r,"x"]
&
G
\end{tikzcd}

\end{center}
is an \(S\)-deviation of \(\varepsilon_G\), hence an element of
\(U(G)\); in the notation of \SGARefLink{sga3:VIIA:1.2}{1.2} it is denoted \(\sigma x\).
Moreover, by \SGARefLink{sga3:VIIA:1.2.1}{1.2.1} the map
\[
  \alpha:x\longmapsto\sigma x
\]
is an isomorphism of \(\mathcal O_S(S)\)-modules from
\(\operatorname{Lie}(G)\) onto the submodule
\(\operatorname{Der}(\varepsilon_G)\subset U(G)\) formed by the
\(S\)-derivations of \(\varepsilon_G\). We show that \(\alpha\) is a
homomorphism of Lie algebras.\footnote{See also \SGARefLink{sga3:ii:subsection:4-11}{Expos\'e II, 4.11}.}

Let
\[
  \rho':U(G)\longrightarrow\operatorname{Dif}_{G/S}
\]
be the algebra homomorphism which sends an \(S\)-deviation \(d\) of
\(\varepsilon_G\) to the left-invariant differential operator
\({}^Gd\), as in \SGARefLink{sga3:VIIA:2.2}{2.2}.

Let
\[
  \rho:G\longrightarrow\operatorname{Aut}_S(G)
\]
be the homomorphism of group functors which sends an \(S\)-morphism
\(g:T\to G\) to right translation by \(g\) on \(G_T\), namely
\begin{center}
  \phantomsection
  \label{sga3:VIIA:2.5-rho-right-translation}
  {\small
\begin{tikzcd}[column sep=normal]
G_T\simeq T\times_TG_T
  \arrow[r,"G_T\times g"]
&
G_T\times_TG_T
  \arrow[r,"m_T"]
&
G_T\,.
\end{tikzcd}}

\end{center}
Recall also from \SGARefLink{sga3:VIIA:1.5}{1.5} and II, \SGARefLink{sga3:ii:subsection:3-14}{3.14} that
\[
 \operatorname{Lie}(\operatorname{Aut}G)
 =
 \operatorname{Lie}(\operatorname{Aut}_S(G)/S)(S)
\]
identifies with the infinitesimal automorphisms of \(G\), that is, the
automorphisms of \(I_S\times G\) which induce the identity on \(G\).

Since \(\rho\) is a monomorphism, so is
\[
 \operatorname{Lie}(\rho):
 \operatorname{Lie}(G/S)
 \longrightarrow
 \operatorname{Lie}(\operatorname{Aut}_S(G)/S)
\]
(see, for example, Expos\'e II, editorial note~50). Consequently
\[
  \operatorname{Lie}(\rho):
  \operatorname{Lie}(G)\longrightarrow
  \operatorname{Lie}(\operatorname{Aut}G)
\]
is injective.

On the other hand, by \SGARefLink{sga3:VIIA:1.5}{1.5} the map \(\beta\), which assigns to every
infinitesimal automorphism \(u\) of \(G\) the differential operator
\(D_u\),
\begin{center}
  \phantomsection
  \label{sga3:VIIA:2.5-beta-sequence}
  {\small
\begin{tikzcd}[column sep=normal]
G\simeq S\times G
  \arrow[r,"\sigma\times G"]
&
I_S\times G
  \arrow[r,"u"]
&
I_S\times G
  \arrow[r,"\operatorname{pr}_2"]
&
G
\end{tikzcd}}

\end{center}
is an isomorphism from
\(\operatorname{Lie}(\operatorname{Aut}G)\) onto the Lie subalgebra
of \(\operatorname{Dif}_{G/S}\) formed by the
\(p^{-1}(\mathcal O_S)\)-derivations of \(\mathcal O_G\).

For every \(x\in\operatorname{Lie}(G)\), the following commutative
square determines the image of \(x\) under \(\operatorname{Lie}(\rho)\):
\begin{center}
  \phantomsection
  \label{sga3:VIIA:2.5-lie-rho-square}
  \begin{tikzcd}[column sep=huge,row sep=large]
I_S\times G
  \arrow[r,"\operatorname{Lie}(\rho)(x)"]
  \arrow[d,"x\times G"']
&
I_S\times G
  \arrow[d,"\operatorname{pr}_2"]
\\
G\times G
  \arrow[r,"m"]
&
G\,.
\end{tikzcd}

\end{center}
It follows from this diagram that the image of
\(\operatorname{Lie}(\rho)(x)\) under \(\beta\) is the composite
deviation
\begin{center}
  \phantomsection
  \label{sga3:VIIA:2.5-composite-deviation}
  {\small
\begin{tikzcd}[column sep=normal]
G\simeq S\times G
  \arrow[r,shift left=.55ex,"\sigma\times G"]
  \arrow[r,shift right=.55ex,"s\times G"']
&
I_S\times G
  \arrow[r,"G\times x"]
&
G\times G
  \arrow[r,"m"]
&
G
\end{tikzcd}}

\end{center}
which, by the editorial note to \SGARefLink{sga3:VIIA:2.2}{2.2}, is precisely
\[
  {}^G(\sigma x)=\rho'(\alpha(x)).
\]
Thus one obtains the commutative diagram
\begin{center}
  \phantomsection
  \label{sga3:VIIA:2.5-lie-algebra-square}
  \begin{tikzcd}[column sep=huge,row sep=large]
\operatorname{Lie}(G)
  \arrow[r,hook,"\operatorname{Lie}(\rho)"]
  \arrow[d,hook,"\alpha"']
&
\operatorname{Lie}(\operatorname{Aut}G)
  \arrow[d,hook,"\beta"]
\\
U(G)
  \arrow[r,hook,"\rho'"]
&
\operatorname{Dif}_{G/S}
\end{tikzcd}

\end{center}
in which \(\operatorname{Lie}(\rho)\), \(\beta\), and \(\rho'\) are
homomorphisms of Lie algebras. Since \(\rho'\) is injective,
\(\alpha\) is also a homomorphism of Lie algebras. Hence:

\medskip
\noindent\textbf{Proposition.}
\phantomsection
\label{sga3:VIIA:2.5-proposition}
\textit{The map \(\alpha\) is an isomorphism of
\(\mathcal O_S(S)\)-Lie algebras from \(\operatorname{Lie}(G)\) onto
the Lie algebra of \(S\)-derivations of \(\varepsilon_G\); the latter
is itself isomorphic, via \(\operatorname{Lie}(\rho)\), to the Lie
algebra of left-invariant \(S\)-derivations of \(G\).}
\footnote{There exist Lie algebras \(\mathfrak g\) over a ring \(A\)
for which the canonical map \(\mathfrak g\to U(\mathfrak g)\) is not
injective; see Bourbaki, \emph{Lie Groups and Lie Algebras}, I.2,
Exercise~9. The result above shows, since \(\alpha\) factors through
\[
  \operatorname{Lie}(G)\longrightarrow
  U(\operatorname{Lie}(G))\longrightarrow U(G),
\]
that this cannot happen for ``algebraic'' Lie algebras, namely those
of the form \(\operatorname{Lie}(G)\) for an \(A\)-group scheme \(G\).}

